\begin{tikzpicture}[x=3.5mm, y=6mm, tr/.style={fill=blue!35, draw=white}, te/.style={fill=orange!70, draw=white}, va/.style={fill=green!45, draw=white}]
\foreach \row/\name/\note in {0/随机划分/{样本独立、同分布时才成立；时序数据中\\相邻时刻会分处训练集与测试集},
  -1.7/时间顺序划分/{用过去预测未来：训练在前，\\验证居中，测试在后},
  -3.4/滚动起点评估/{测试窗口逐步后移，\\得到误差的分布而不是一个数},
  -5.1/分组或空间留出/{同一路段、同一起事故或同一区域\\只出现在训练集或测试集一侧}}
{
  \node[anchor=south west, font=\footnotesize\bfseries, inner sep=1pt] at (0,\row+0.38) {\name};
  \node[anchor=west, tag, align=left] at (25,\row) {\note};
}
\foreach \i in {0,...,23} {
  \pgfmathtruncatemacro{\m}{mod(\i*7+3,10)}
  \ifnum\m<2 \fill[te] (\i,-0.35) rectangle (\i+1,0.35); \else \ifnum\m<3 \fill[va] (\i,-0.35) rectangle (\i+1,0.35); \else \fill[tr] (\i,-0.35) rectangle (\i+1,0.35); \fi\fi
}
\fill[tr] (0,-2.05) rectangle (16,-1.35); \fill[va] (16,-2.05) rectangle (20,-1.35); \fill[te] (20,-2.05) rectangle (24,-1.35);
\foreach \k/\y in {0/-3.25,1/-3.55,2/-3.85} {
  \fill[tr] (0,\y-0.12) rectangle ({14+3*\k},\y+0.12); \fill[te] ({14+3*\k},\y-0.12) rectangle ({17+3*\k},\y+0.12);
}
\foreach \i in {0,...,23} {
  \pgfmathtruncatemacro{\g}{mod(\i,6)}
  \ifnum\g=4 \fill[te] (\i,-5.45) rectangle (\i+1,-4.75); \else \fill[tr] (\i,-5.45) rectangle (\i+1,-4.75); \fi
}
\draw[-Stealth, black!50] (0,-6.0) -- (24.3,-6.0) node[right, tag]{时间};
\fill[tr] (2,-6.8) rectangle (3,-6.4); \node[tag, anchor=west] at (3.2,-6.6) {训练};
\fill[va] (7,-6.8) rectangle (8,-6.4); \node[tag, anchor=west] at (8.2,-6.6) {验证};
\fill[te] (12,-6.8) rectangle (13,-6.4); \node[tag, anchor=west] at (13.2,-6.6) {测试};
\end{tikzpicture}
